\section{Malware Classification}

\subsection{Classification for Attack Detection}
Unlike anomaly detection, classification of malware is a supervised learning problem and requires representative data from \it{all} classes. We generally assume that \b{unknown attacks are related to known attacks}.
% \b{Note}: All practical detection approaches are approximations!

\subsubsection{Data Collection}
The required data can be collected using either \b{Honeypots} (active or passive baits) or \b{Forensic Analysis} (analysis of past security incidents, quarantined malware, physhing mails, ...).\\[1em]
\b{Security Risks}: Systems can be undermined by \b{Poisoning} (injecting bad data during training), \b{Mimicry} (making malware look benign), or \b{Red Herrings} (flooding the system with false alarms).

\subsubsection{Feature Extraction}
\definition{Static Analysis}{Extracting features by analyzing the malware executable file without executing it. Features include strings, byte \f{n}-grams, and Portable Executable (PE) headers.
\begin{itemize}
\item \b{Pros}: Extremely fast, inherently safe (no execution risk), and can theoretically analyze all execution paths to achieve full code coverage.
\item \b{Cons}: Highly vulnerable to static evasion techniques like packing, obfuscation, and encryption, which hide the true malicious payload.
\end{itemize}}

\definition{Dynamic Analysis}{Extracting features by running the malware in a controlled, monitored environment (sandbox) and recording its runtime behavior. Features include API calls, file system changes, and network traffic.
\begin{itemize}
\item \b{Pros}: Resilient against packing and static obfuscation, as the malware is forced to unpack itself in memory to execute.
\item \b{Cons}: Computationally expensive (slow), may not trigger all execution paths (low code coverage), and is susceptible to environment evasion (e.g., malware detecting the sandbox and remaining dormant).
\end{itemize}}


\subsection{Models for Malware Classification}
\begin{itemize}
    \item \b{Perceptron}: A weighted model that learns by iteratively updating its weights \f{w}. In the single perceptron case, \f{w} defines the orientation of the hyperplane, and it only converges for linearly separable data. Decision function:
    \cf{f(z) = sign(\langle\phi(z),w\rangle)}
    \item \b{Two-Class SVM}: A model with weight \f{w} and bias \f{b} (defining the hyperplane) that is optimized such that the margin \f{M=\frac{2}{||W||}} between the hyperplane and both classes is maximized, subject to correctly separating the classes:
    \cf{\min_{w,b}||w||\quad\text{s.t.}\quad y_i\cdot\left(\langle\phi(x_i),w\rangle+b\right)\geq 1 - \xi_i}
    Just as before, to handle non-linear data (i.e. outliers), slack variables ($\xi_i \geq 0$) allow some samples to fall on the wrong side of the margin. Furthermore, we can use the Kernel trick to separate data in higher-dimensional spaces.
\end{itemize}





